\documentclass[UTF8]{ctexart}
\usepackage{geometry}
\usepackage{hyperref}
\usepackage{listings}
\usepackage{xcolor}

\geometry{a4paper, left=2.5cm, right=2.5cm, top=2.5cm, bottom=2.5cm}

\title{Agent开发周报 260208}
\author{刘德辰}
\date{2026年02月08日}

\begin{document}

\maketitle

\section{规则驱动智能体（Rule-Driven Agent）核心功能实现}

本周重点完成了规则驱动智能体的核心架构与功能实现，该系统允许客户通过自然语言配置回复规则，系统在运行时根据 Router 的匹配结果自动按规则生成稳定且自然的回复。

\subsection{架构设计说明}

本系统的架构设计采用了分层设计思路：
\begin{itemize}
    \item \textbf{总体架构}：基于 RAG（Retrieval-Augmented Generation）通用方案，采用向量检索 + LLM 生成的经典架构模式。
    \item \textbf{Config 架构}：针对规则配置场景，设计了"双轨交付"和"对话式创建"架构（Ask / Build / Reply 三类技能），该设计来源于对用户需求和使用场景的深入分析，通过对话式配置降低客户使用门槛，通过编译机制规范规则格式，保证规则的可执行性和安全性，是结合了业务场景的创新设计。
\end{itemize}

\subsection{Config 架构与交互模式设计（核心亮点）}

\subsubsection{对话式创建：降低配置门槛的核心创新}

系统的核心亮点是\textbf{对话式创建}机制，让客户用"人话"配置规则，无需理解技术概念。

\begin{itemize}
    \item \textbf{Rule Asker（规则配置追问向导）}：采用"需求访谈"式的对话方式，像与产品经理沟通一样自然
    \begin{itemize}
        \item 通过自然语言逐步追问收集规则信息（规则名称、触发描述、回复目标、关键要点、必填信息、模板、禁区、语调等）
        \item 每轮最多追问 1-2 个问题，避免人脑信息过载，降低配置门槛
        \item 当信息齐全时输出确认摘要，引导用户确认保存
        \item 支流式输出（SSE），提供实时的交互反馈
        \item 产物为 \texttt{rule\_draft}（草稿结构），采用业务语言表达，存储在 \texttt{rule\_drafts} 表中
    \end{itemize}
    
    \item \textbf{Rule Builder（规则编译器）}：分析草稿和 asker 对话全文，将其编译为规范化规则格式
    \begin{itemize}
        \item 补全默认值、规范化字段、生成规则 ID、注入安全底线
        \item 使用 embedding service 生成向量，用于后续规则匹配
        \item 输出严格的 JSON 格式，可校验、可存储、可版本化
        \item 保存前自动进行冲突检测，防止规则重复或语义冲突
    \end{itemize}
\end{itemize}

\subsubsection{双轨交付：交互与编译的物理隔离}

系统延续"双轨"理念，将\textbf{交互式对话}与\textbf{结构化编译}物理隔离，实现配置体验与系统可靠性的最佳平衡：
\begin{itemize}
    \item \textbf{交互式轨道}：Rule Asker 和 Rule Reply 采用 Markdown 流式输出，提供流畅的用户体验
    \item \textbf{结构化轨道}：Rule Builder 输出 Strict JSON，经过兜底处理和 check 后保存，保证规则格式的规范性和执行稳定性
\end{itemize}

\subsection{Embedding Service 层（规则匹配服务）}

\begin{itemize}
    \item \textbf{独立匹配服务}：新增了 \texttt{match\_rules} 工具，作为 Router Pipeline 中的独立服务层，专门负责规则的向量化与相似度匹配。该服务不依赖 Router LLM 的判断，使用专门的 Embedding 模型进行向量化计算。
    \item \textbf{混合检索策略}：采用 \textbf{Anchor + Examples} 的混合检索策略，通过 \texttt{match\_text}（锚点）保证召回率，通过 \texttt{examples}（用户真实语料示例）提升精确率和排序质量。融合算法使用 Max Fusion 策略，取锚点得分和示例最大得分的较高值作为最终得分。
    \item \textbf{向量化实现}：使用阿里云 QWEN \texttt{text-embedding-v1}模型进行向量化，支持规则保存时自动生成 embedding。向量数据存储在 \texttt{rules.embedding} 字段中，采用 JSON 格式存储锚点和示例的向量数组。
    \item \textbf{当前实现方式}：受限于当前测试 demo 环境后端部署在 Cloudflare Workers 和 D1 数据库上，当前实现采用\textbf{全量加载 + 内存计算}的方式：
    \begin{itemize}
        \item 从 D1 数据库全量加载所有启用规则（当前限制最多 500 条）
        \item 在 Workers 内存中对每个规则计算余弦相似度
        \item 按相似度排序并返回 TopK 结果
        \item \textbf{限制}：D1（SQLite）不支持向量索引，无法进行高效的向量检索；Workers 有执行时间限制（通常 30 秒），当规则数量增长时性能会急剧下降
    \end{itemize}
    \item \textbf{向量数据库未实现}：由于上述环境限制，当前\textbf{未实现向量数据库}（如 Milvus、Qdrant、Pinecone 等），向量检索相关的高效索引和检索功能尚未实现。这是当前架构的主要技术债务，生产环境中会引入专门的向量数据库来解决该问题。
    \item \textbf{TopK 检索}：实现了相似度计算和排序算法，支持自定义 \texttt{top\_k}（默认 5）和 \texttt{min\_score}（默认 0.3）参数，返回匹配度最高的规则列表。
\end{itemize}

\subsection{Agent 技能（Ask / Build / Reply）}

系统新增三个技能，严格遵循"双轨交付"理念，实现了规则配置、编译和执行的完整闭环。

\begin{itemize}
    \item \textbf{Rule Asker（规则配置追问向导）}：面向客户的对话式配置 Agent，通过自然语言逐步追问收集规则信息。每轮最多追问 1-2 个问题，收集包括规则名称、触发描述、回复目标、关键要点、必填信息、模板、禁区、语调等字段。当信息齐全时输出确认摘要，引导用户确认保存。产物为 \texttt{rule\_draft}（草稿结构），存储在 \texttt{rule\_drafts} 表中。
    \item \textbf{Rule Builder（规则编译器）}：面向系统的结构化编译器 Agent，负责将 \texttt{rule\_draft} 编译为 \texttt{rule\_json}。职责包括补全默认值、规范化字段、生成规则 ID、注入安全底线、生成 embedding 向量。输出为严格的 JSON 格式，可校验、可存储、可版本化。
    \item \textbf{Rule Reply（规则执行器）}：面向运行时的规则执行工具，Router 已提供命中规则 IDs 和匹配分数。负责读取规则详情，按规则模板生成最终回复（Markdown 流式输出）。支持模板填空、缺参追问、安全底线检查等功能。
\end{itemize}

\subsection{Router 集成与决策流程}

\begin{itemize}
    \item \textbf{Router integration}: Runtime ALWAYS runs \texttt{match\_rules} before Router and injects results; Router does not call the tool.
    \item \textbf{Decision logic}: Router decides based on precomputed \texttt{match\_rules} results (top1 > 0.7 => prefer rule\_reply, even for emergency).
    \begin{itemize}
        \item 如果 top1 匹配分数 > 0.7 → 派单 \texttt{rule\_reply}，传入 \texttt{hit\_rules} 和 \texttt{user\_query}
        \item 如果匹配分数较低或无匹配 → 继续走其他 Worker（如 \texttt{consult\_omni}）或拒绝
    \end{itemize}
    \item \textbf{场景表兼容}：规则匹配机制与现有的工作场景表机制兼容，Router 可以结合场景表和规则匹配结果做出最终决策。
    \item \textbf{决定因素}：任务分发仍由 Router agent 最终决定，匹配结果为参考因素。
\end{itemize}

\section{规则冲突检测机制}

实现了规则保存前的冲突检测功能，防止规则重复或语义冲突。

\begin{itemize}
    \item \textbf{自检索机制}：在 \texttt{rule\_builder} 保存规则前，对当前草稿的 \texttt{match\_text} 生成 Embedding 向量，并在规则库中执行 Top-K 检索（排除当前规则 ID）。
    \item \textbf{三级阈值判定}：
    \begin{itemize}
        \item \textbf{红区 (> 0.92)}：高度冲突，视为重复规则，阻断保存，提示用户直接修改旧规则
        \item \textbf{黄区 (0.80 - 0.92)}：潜在冲突，语义相近，提示用户建议合并，询问是否强制保存
        \item \textbf{绿区 (< 0.80)}：安全，直接保存
    \end{itemize}
    \item \textbf{冲突检测工具}：实现了 \texttt{detectRuleConflict} 函数，支持在规则新建和编辑时进行冲突检测，确保规则库的质量和一致性。
\end{itemize}

\section{数据库架构扩展}

新增了规则系统的数据库表结构，支持规则和草稿的存储与管理。

\begin{itemize}
    \item \textbf{rules 表}：存储编译后的可执行规则，包含 \texttt{id}、\texttt{name}、\texttt{match\_text}、\texttt{enabled}、\texttt{priority}、\texttt{version}、\texttt{embedding}、\texttt{data}（规则 JSON 数据）等字段。建立了 \texttt{enabled} 和 \texttt{priority} 索引，支持规则的启用/禁用和优先级管理。
    \item \textbf{rule\_drafts 表}：存储规则配置过程中的草稿数据，包含 \texttt{session\_id}（主键）、\texttt{status}、\texttt{mode}（create/edit）、\texttt{rule\_id}、\texttt{draft}（草稿 JSON 数据）等字段。支持多会话并发配置规则。
\end{itemize}

\section{规则数据结构设计}

设计了 \texttt{rule\_draft}（草稿）和 \texttt{rule\_json}（可执行规则）两阶段数据模型。\texttt{rule\_draft} 采用业务语言字段（\texttt{name}、\texttt{match\_text}、\texttt{examples}、\texttt{reply\_goal}、\texttt{key\_points}、\texttt{required\_info}、\texttt{template}、\texttt{safe\_defaults}、\texttt{do\_not\_say}、\texttt{tone} 等），降低配置门槛；\texttt{rule\_json} 为系统内部严格 JSON 格式，包含完整执行逻辑和 \texttt{embedding} 字段。

\section{系统架构优化}

延续 Router 与 Worker 职责分离架构，规则工具通过统一接口接入系统，\texttt{rule\_asker} 和 \texttt{rule\_reply} 支持流式输出（SSE）。

\section{文档与规范}

完成《规则列表系统开发文档.md》和三类技能的 SKILL.md 文档，在关键代码中添加详细注释。

\section{技术债务与限制}

\textbf{向量数据库未实现}：受限于 Cloudflare Workers + D1 环境，当前采用全量加载 + 内存计算方式，无法支持大规模规则库的高效检索。D1 不支持向量索引，Workers 有执行时间和内存限制，当规则数量超过数百条时性能会急剧下降。需引入专门的向量数据库解决方案。

\section{下周计划}

\begin{enumerate}
    \item \textbf{向量数据库集成}：调研并实现向量数据库方案，重构 \texttt{match\_rules} 和 \texttt{detectRuleConflict} 使用向量数据库进行高效检索。
    \item \textbf{规则测试功能}：实现规则测试工具，支持保存前测试匹配效果和回复质量。
    \item \textbf{规则版本管理}：完善版本控制机制，支持历史版本查看和回滚。
    \item \textbf{规则模板库}：建立常用规则模板库，降低配置成本。
\end{enumerate}

\end{document}